\input{preamble}

\begin{document}

% Remove page number
\clearpage
\thispagestyle{empty}

\begin{bfseries}
\begin{center}

\includegraphics[scale=0.45]{ufc.png} \\
\vspace{-4pt} 
UNIVERSIDADE FEDERAL DO CEARÁ \\
\vspace{4pt} 
CENTRO DE CIÊNCIAS \\
\vspace{4pt} 
DEPARTAMENTO DE FÍSICA \\
\vspace{4pt}
LABORATÓRIO DE FÍSICA EXPERIMENTAL PARA ENGENHARIA \\
\vspace{4pt}
SEMESTRE 2020.1 \\


\vspace*{\fill}
\textbf{PRÁTICA XX - TÍTULO}
\vspace*{\fill}

\end{center}

\begin{itemize}[leftmargin=*]
    \setlength{\itemsep}{0pt}
    \item[] ALUNO: XXXXXXXXXXXXXXXXXX
    \item[] MATRÍCULA: XXXXXXXXXXXXXXX
    \item[] CURSO: XXXXXXXXXXXXXXXXXX
    \item[] TURMA: XXXXXXXXXXXXXXXXXX
    \item[] PROFESSOR: XXXXXXXXXXXXXX
    \item[] DATA E HORA DA REALIZAÇÃO DA PRÁTICAS 00/00/0000 ÀS 00:00h
\end{itemize}

\end{bfseries}
\newpage

\section{OBJETIVOS}
Citar os objetivos da prática em questão. Aqui é permitido reproduzir os objetivos citados no Roteiro de Práticas.

\section{MATERIAL}
Listar o material utilizado na prática. Aqui é permitido reproduzir a lista de material citada no Roteiro de Práticas.

\section{INTRODUÇÃO}
Início da parte textual do trabalho. Tem como finalidade dar ao leitor uma visão concisa do tema da prática, redigindo um texto de sua própria autoria sobre a fundamentação teórica envolvida na prática em questão. Aqui se pretende que o aluno exercite a habilidade de escrever em linguagem científica. Para redigir a introdução recomenda-se que o aluno pesquise em manuais, livros-texto e/ou Internet.

Todos os autores citados devem ter a referencia incluída em lista no final no trabalho. Não será tolerada a cópia total ou parcial sem a devida referência. A reprodução de figuras obtidas de outras fontes é permitida desde que seja citada a fonte, devidamente, no texto e nas referências.

Todo texto deve ser digitado em fonte Times New Roman, tamanho 12, inclusive a capa, com exceção das citações com mais de três linhas, notas de rodapé, paginação, legendas e fontes das ilustrações e das tabelas, que devem ser em fonte Times New Roman tamanho 10 e espaçamento entre linhas de 1,5. Com Margens: Superior e esquerda em três centímetros; inferior e direita em dois centímetros. O texto deve ser justificado, exceto as referências, no final do trabalho, que devem ser alinhadas a esquerda.

\jump

A seguir, têm-se os exemplos de como fazer citações e referências.

\jump

Para Siss (2012) as políticas de ação afirmativas constituem políticas públicas, estatais e de caráter compulsório, elaboradas e implementadas pelo Estado, ou seja, é o Estado em ação.

Segundo Bastos e Keller (2006, p. 38), “A leitura é um processo que envolve algumas habilidades, entre as quais a interpretação do texto e a sua compreensão.”

As organizações testemunharam uma redução da validade de seu conhecimento durante este período e começaram a perceber que já não era possível confiar em Instituições de Ensino Superior para desenvolver a sua mão de obra (TARAPANOFF, 2006).

O discurso jurídico, que hoje se apresenta com um novo perfil, dispõe de um acervo variado de opções para ser construído, pois, “[...] agrega valores, impõe condutas, conduz instituições, movimenta riquezas, opta por visões de mundo e, portanto, sustenta uma ideologia.” (BITTAR, 2001, p. 181).

\begin{leftquote}[]
A citação direta com mais de três linhas deve ser destacada com recuo de 4 cm da margem esquerda, em letra menor do que a utilizada no texto (tamanho 10), sem as aspas e com espaçamento simples entrelinhas. A citação deverá ser separada do texto que a precede e a sucede [...] (UNIVERSIDADE FEDERAL DO CEARÁ, 2013, p. 96).
\end{leftquote}

\begin{leftquote}[De acordo com Mueller e Perucchi (2014, p. 16)]
A expressão ciência e tecnologia e inovação, comumente abreviada para CT\&I, reflete a interação e a interdependência entre essas áreas. Pela sua natureza, o conhecimento científico, tecnológico e a inovação dependem da comunicação, e a comunicação desses conhecimentos é objeto de estudo da ciência da informação.
\end{leftquote}

\jump

As ilustrações (fotografias, gráficos, mapas, plantas, quadros) e tabelas devem ser citados e inseridos o mais próximo possível do trecho a que se referem.

\jump

% http://pgfplots.net/tikz/examples/bar-plot/ (the "right" way to do it)
\begin{container}{10cm}{Gráfico 1 – Distribuição dos documentos analisados por programa de pós-graduação}{Fonte: Elaborado pelo autor.}

\includegraphics[width=10cm]{graph.png}

\end{container}

\begin{container}{10cm}{Figura 1 – Organização do conhecimento/Representação do
conhecimento, Organização da informação/Representação da informação}{Fonte: Lara e Smit (2010, com adaptações).}

\includegraphics[width=10cm]{image.png}

\end{container}

\jump

\begin{container}{10cm}{Tabela 1 – Distribuição dos documentos analisados por programa de pós-graduação}{Fonte: elaborada pelo autor.}
\begin{center}
\begin{tabular}{ | c | c | c | c | }
    \hline
    Programas de pós-graduação & Teses & Dissertações & Total\\
    \hline
    \hline
    Cirurgia & 1 & 1 & 2 \\
    \hline
    Enfermagem & 4 & 4 & 8 \\
    \hline
    Engenharia Civil & 2 & 8 & 10 \\
    \hline
    Farmacologia & 8 & 6 & 14 \\
    \hline
    Física & 3 & 6 & 9 \\
    \hline
    Química Inorgânica & 4 & 1 & 5 \\
    \hline
    Total & 22 & 26 & 48 \\
    \hline
\end{tabular} 
\end{center}
\end{container}

\section{PROCEDIMENTO}
Nesta parte o aluno deve descrever detalhadamente os procedimentos realizados (em primeira pessoa); apresentar os resultados das medidas (tabelas, gráficos, cálculos, etc). Atenção: anotar sempre as unidades.


\section{QUESTIONÁRIO}
No questionário devem constar as perguntas e as respostas de acordo com o roteiro de prática do ano em curso. 


\section{CONCLUSÃO}
Parte final do texto na qual se apresentam as conclusões apoiadas no desenvolvimento do assunto. É a recapitulação sintética dos resultados obtidos. 

Na conclusão o aluno deve apresentar uma discussão sobre os resultados obtidos em função dos objetivos do experimento. Cabe também comentar sobre as possíveis fontes de erros.

\newpage

\begin{center}
    \textbf{REFERÊNCIAS}
\end{center}

\jump

{\parindent0pt % disables indentation for all the text between { and }

BASTOS, Cleverson Leite; KELLER, Vicente. \textbf{Aprendendo a aprender}: introdução à metodologia científica. 19. ed. Petrópolis: Vozes, 2006.

\jump

BITTAR, Eduardo Carlos Bianca. \textbf{Linguagem jurídica.} São Paulo: Saraiva, 2001.

\jump

LARA, Marilda Lopes Ginez de; SMIT, Johanna Wilhelmina. \textbf{Temas de pesquisa em Ciência da Informação no Brasil.} São Paulo: Escola de Comunicações e Artes da Universidade de São Paulo, 2010. Disponível em:
http://www.repositoriobib.ufc.br/000005/00000588.pdf. Acesso em: 21 jan. 2012.

\jump

MUELLER, Suzana Pinheiro Machado; PERUCCHI, Valmira. Universidades e a produção de patentes: tópicos de interesse para o estudioso da informação tecnológica. \textbf{Perspectivas em Ciência da Informação}, Belo Horizonte, v. 19, n. 2, p. 15-36, 2014.

\jump

SISS, Ahyas. Afro-brasileiros e Educação Superior: notas para debates. In: COSTA, Hilton; PINHEL, André; SILVEIRA, Marcos Silva da (org.). \textbf{Uma década de políticas afirmativas}: panorama, argumentos e resultados. Ponta Grossa: Editora UEPG, 2012. p. 18-26.

\jump

TARAPANOFF, K. Educação corporativa. In: CONGRESSO IBEROAMERICANO DE
GESTÃO DO CONHECIMENTO E INTELIGÊNCIA COMPETITIVA, 1., 2006, Curitiba.
\textbf{Anais} [...]. Curitiba: CIETEP, 2006. Disponível em: http://www.gecic.com.br. Acesso em: 22 out. 2006. p. 59-70.

\jump

UNIVERSIDADE FEDERAL DO CEARÁ. Biblioteca Universitária. \textbf{Guia de normalização de trabalhos acadêmicos da Universidade Federal do Ceará}. Fortaleza, 2013.

}

\end{document}
